\chapter{Layout: flexbox in detail}
\label{ch:flexbox2}

\begin{aside}
Flexbox is the layout engine behind nearly all modern UI.
\maya{}'s implementation is small but covers the cases that
matter for terminals.
\end{aside}

\section{Concepts}

\begin{itemize}[leftmargin=*]
  \item Main axis vs cross axis.
  \item Flex grow, flex shrink, flex basis.
  \item Justify content, align items.
\end{itemize}

\section{Main axis}

In a row, main is horizontal; in a column, main is
vertical. Children flow along the main axis.

\section{Flex factors}

\begin{lstlisting}
auto row = row(
    sidebar() | flex(0),    // fixed
    content() | flex(1),    // takes the rest
    inspector() | flex(0)   // fixed
);
\end{lstlisting}

\code{flex(0)} means ``don't grow''. \code{flex(1)} means
``share the leftover space''.

\section{Distributing}

If three children all have \code{flex(1)}, the leftover
splits in three. If they're \code{flex(1)}, \code{flex(2)},
\code{flex(1)}, the middle gets twice as much.

\section{Shrink}

When content overflows, shrink redistributes:
\code{shrink(0)} forbids shrinking; \code{shrink(1)} allows.

\section{Justify content}

How children align on the main axis when there's leftover
space:

\begin{itemize}[leftmargin=*]
  \item start: pack to start.
  \item end: pack to end.
  \item center: pack to centre.
  \item space-between: gaps between children.
  \item space-around: gaps including ends.
  \item space-evenly: equal gaps.
\end{itemize}

\section{Align items}

How children align on the cross axis:
start, end, center, stretch (default).

\section{Wrap}

\code{flex\_wrap} flows children onto multiple lines when
they don't fit. \maya{} supports this but rarely uses it
for TUI layouts.

\section{Recap}

\begin{recap}
\begin{itemize}[leftmargin=*]
  \item Main axis flow + cross axis alignment.
  \item Flex factors share leftover space.
  \item Shrink controls overflow behaviour.
  \item Justify and align handle gaps.
\end{itemize}
\end{recap}

\section*{Workshop}

\begin{enumerate}[leftmargin=*]
  \item Build a three-pane layout with flex factors.
  \item Add justify-space-between for a footer.
  \item Centre a dialog with align-center.
\end{enumerate}
